\documentclass[11pt,a4paper]{article}
\usepackage[margin=25mm]{geometry}
\usepackage{amsmath,amssymb,amsthm,hyperref}
\setlength{\emergencystretch}{2em}
\newtheorem{lemma}{Lemma}
\newtheorem{theorem}[lemma]{Theorem}
\newtheorem{corollary}[lemma]{Corollary}
\newcommand{\gen}[1]{\langle#1\rangle}
\title{Localized eligible edges and the final internal orbit cases}
\author{Human Compression project: L044 continuation}
\date{}
\begin{document}\maketitle
Throughout, $G=\gen{\rho_0,\ldots,\rho_{r-1}}$ is transitive
on $\Omega$, and its involutory generators satisfy the string
commutations. An edge of colour $j$ is \emph{ambient-eligible}
when its endpoints belong to different orbits of
$G_j=\gen{\rho_i:i\ne j}$. All interval results below use these
actual ambient omissions. A local omission condition is not substituted.

\begin{lemma}[Localization includes fixed vertices]\label{local}
Let $I$ be a nonempty contiguous interval of colours and put $T=G_I$.
Suppose $V$ is a proper $T$-invariant subset of $\Omega$, every colour
in $I$ has at least two ambient-eligible edges, and \emph{every}
ambient-eligible edge of these colours has both endpoints in $V$.
Then, counting all actual $T$-orbits in $V$, including singletons,
\[
                     2|I|\le |V|-o_T(V).                 \tag{1}
\]
\end{lemma}
\begin{proof}
Use the complete accepted ambient-interval theorem and its
minimum-square normal form. Every selected edge lies in $V$, in
every eligible selection, by the stated confinement hypothesis.
Thus the minimization and its exchanges cannot move an edge outside
$V$. Each nontrivial $T$-orbit in $V$ is a proper ambient orbit,
so the theorem supplies a tree component in it. Its vertices minus
selected edges are at least one. Every $T$-fixed point in $V$ is
isolated and supplies the same unit. Summing over all orbits in $V$
proves (1). Merely choosing an initial selection in $V$ would not
justify this argument: the confinement of \emph{all} eligible edges
is essential.
\end{proof}

\begin{corollary}[The odd-anchor budget is impossible]\label{anchor}
Retain the entire six-point head and normalized critical assignment
of \emph{Strict row exits and critical internal orbit budgets},
including its actual full IP and natural alternating tail.
Then its remaining first-assigned internal odd-orbit configuration
is impossible.
\end{corollary}
\begin{proof}
The accepted reduction leaves exactly seven $H=G_{\ge5}$-fixed
points, a first-assigned internal orbit $O$ of size $2q$ with odd
$q\ge3$, and an odd orbit $V$ of size $m\ge5$. The crossing
involution $\tau=\rho_3$ is free on $O$ and fixes $V$.
The assigned labels on $O$ are $F_O=\{5,\ldots,q+3\}$;
their pair images form the ordinary path, and every other tail
pair image on $O$ is trivial. All remaining labels form the
terminal interval $I=F_V$, of size $(m-1)/2$.

For $j\in I$, every possible $j$-edge on $O$ joins a
$\tau$-pair. It is ineligible, since $\tau\in G_j$.
The seven head points are fixed. Therefore all ambient-eligible
edges of the colours in $I$ lie in $V$, and each colour has two
by its assignment. With $T=G_I$, the actual local
$K=G_{>5}$ is intransitive on $V$ by the normalized assignment,
and $T\le K$. Hence $o_T(V)\ge2$. Lemma~\ref{local} gives
\[
                     m-1=2|I|\le m-o_T(V)\le m-2,
\]
a contradiction. No assertion about a quotient intersection
property or signed lifts is used.
\end{proof}

\begin{lemma}[A two-point remainder at the exceptional head]\label{pair}
Suppose $n=2r\ge18$ and there is a partition
\[
       \Omega=D\mathbin{\dot\cup}\{b\}
                    \mathbin{\dot\cup}E\mathbin{\dot\cup}V,
       \qquad D=\{1,\ldots,6\},\quad E=\{u,v\}.
\]
Assume the following actual marked-group hypotheses:
\begin{itemize}
\item $\gen{\rho_0,\rho_1,\rho_2}$ is transitive and supported
on $D$; $\tau=\rho_3=(6\ b)(u\ v)$; $\sigma=\rho_4$ fixes $D$.
\item Full ambient IP holds and $G_{\ge3}=A_Z$ naturally,
where $Z=\Omega\setminus\{1,\ldots,5\}$.
\item $H=G_{\ge5}$ fixes $D\cup\{b\}$, preserves $E$,
and is transitive on $V$. The group
$B=\gen{\sigma,H}$ is transitive on $\Omega\setminus D$.
\item Every colour greater than $5$ has two ambient-eligible edges.
\end{itemize}
These hypotheses are contradictory. The action of $H$ on $E$
may be transitive or pointwise trivial.
\end{lemma}
\begin{proof}
Put $x=b^\sigma$. Transitivity of $B$ gives $x\ne b$.
The accepted exceptional-head lemma gives $x^\tau=x$.
Since $\sigma$ fixes $D$ and $\tau$ exchanges the two points
of $E$, it follows that $x\in V$. Write $K=G_{>5}$.
String commutation implies that $K$ centralizes $\tau,\sigma$.
It fixes $b$ and hence $x$, and preserves $E$ and $E^\sigma$.

We recall the exact pure-group step of the head lemma.
Conjugating $\tau$ by the transitive group on $D$ and multiplying
two conjugates cancels $(u\ v)$ and produces the pure three-cycles
on $\{b,d,d'\}$. Thus $G_{\le3}$ contains natural
$A_{D\cup\{b\}}$. Its conjugate by $\sigma$, together with
it, gives natural $A_{D\cup\{b,x\}}\le G_{\le4}$.
In particular a pure three-cycle on $\{6,b,x\}$ lies in
\[
             G_{\le4}\cap G_{\ge3}=\gen{\tau,\sigma}.  \tag{2}
\]
Both containment in the natural alternating tail and the equality
in (2) concern actual subgroups, not projected markings.

The involution $\sigma$ cannot preserve $E$, since $H$ does so
and $B$ is transitive. If it moves exactly one point of $E$
outside, rename that point $u$ and put $w=u^\sigma$.
Then $w\in V\setminus\{x\}$ and $\sigma$ fixes $v$.
On the invariant sets $\{6,b,x\}$ and $\{u,v,w\}$,
the rotation $\tau\sigma$ has order three on each:
its two factors are respectively $(6\ b),(b\ x)$ and
$(u\ v),(u\ w)$. Every odd-order element of a group generated
by two involutions is a power of its rotation. A power trivial on
$\{u,v,w\}$ is divisible by three and therefore trivial on
$\{6,b,x\}$. This contradicts the pure three-cycle in (2).

Consequently both points move, say $u^\sigma=w$ and
$v^\sigma=z$, where $x,w,z$ are distinct points of $V$.
The set $Y=V\setminus\{x,w,z\}$ is $K$-invariant:
$K$ preserves $V$ and $\{w,z\}=E^\sigma$ and fixes $x$.
For a colour $j>5$, every edge outside $Y$ is either on $E$
or on $\{w,z\}$, since $D,b,x$ are fixed by $K$.
An edge on $E$ is ineligible because $\tau\in G_j$;
one on $\{w,z\}$ is ineligible because
$\tau^\sigma=(6\ x)(w\ z)\in G_j$.
Thus all eligible edges for the nonempty interval
$I=\{6,\ldots,r-1\}$ lie in $Y$.
But $|V|=n-9$ gives $|Y|=n-12=2(r-6)=2|I|$.
Lemma~\ref{local} implies $2|I|\le |Y|-o_K(Y)<|Y|$,
the final contradiction. All singleton $K$-orbits in $Y$ count
in $o_K(Y)$; none is discarded.
\end{proof}

\begin{corollary}[Empty internal pairs and exchanged fixed points]
Retain the complete normalized critical hypotheses for the two
source heads in \emph{Strict row exits and critical internal orbit
budgets}. An internally swapped actual two-point $H$-orbit is
impossible in both heads. In the exceptional six-point head,
the crossing involution cannot exchange two additional $H$-fixed
points outside the specified seven-point head either.
\end{corollary}
\begin{proof}
Suppose first that $E$ is an actual two-point orbit internally
swapped by $\tau$. No tail edge on $E$ is ambient-eligible.
Let $f$ be the number of $H$-fixed points and $c$ the number
of its nontrivial orbits. All eligible tail edges lie in the
union $W$ of the other nontrivial orbits. Applying
Lemma~\ref{local} to the full tail interval yields
\[
  2(r-k)\le n-f-2-(c-1),\qquad f+c\le2k-1.             \tag{3}
\]
There is another nontrivial orbit, because each tail colour has
an eligible pair whereas $E$ has none. For the singleton head,
$k=2$, $f\ge2$, $c\ge2$ contradict (3).
For the exceptional head $k=5$, $f\ge7$ leaves only $f=7,c=2$.
The other orbit $V$ has odd size $n-9$.
An involution centralizing a transitive action on an odd orbit
fixes it pointwise; it cannot exchange $V$ with $E$.
The literal exceptional head therefore gives
$\tau=(6\ b)(u\ v)$, and every premise of Lemma~\ref{pair}
is exactly the inherited source-head interface. This is impossible.

Finally suppose $\tau$ exchanges two additional $H$-fixed points.
The full ambient-interval theorem for $H$ gives
$2(r-5)\le n-f-c$, hence $f+c\le10$.
Here $f\ge9$ and $c\ge1$, so $f=9,c=1$.
The two additional fixed points form $E$; the sole nontrivial
orbit $V$ has odd size $n-9$ and is fixed pointwise by $\tau$.
Again the literal head gives $\tau=(6\ b)(u\ v)$.
Lemma~\ref{pair} applies with $H$ fixing $E$ pointwise.
This covers the singleton orbits explicitly rather than treating
them as part of the support of $H$.
\end{proof}

\paragraph{Source, dependencies and limits.}
These are bounded continuations of the orbit accounting in
Cameron--Fernandes--Leemans--Mixer, \emph{Highest rank of a polytope
for $A_n$}, arXiv:1605.09173v1, Proposition 5.31, physical
pages 42--43. The complete accepted ambient-interval theorem,
strict internal reductions and critical head interfaces remain
paid; in particular their smaller-degree alternating induction,
full IP and actual supported head requirements are retained.
The confinement lemma and the final two-point rotation obstruction
are the present arguments. Root checked the localized-interval
idea and specifically requested that isolated vertices remain in
the orbit count. The older global middle move retains its separate
authorship and distinct review through the paid interval proof.

The source's support notation alone does not account for exchanged
fixed singleton points; the last corollary handles that possibility
under the explicit critical hypotheses. No unrestricted source
assignment existence is inferred, no local eligibility is promoted
to ambient eligibility, and no signed lift classification is used.
Other source branches, rigidity diagrams, tables and complete L044
recovery remain separate. No new computation, formal theorem,
whole-original credit or question claim is made.
\end{document}
